\section*{Capítulo \ref{ch:b3:exploration-seismology} --- Sismologia de exploração}

\begin{solution}{exo:b3:exploration-seismology:1}
$Z = \rho v$: giz $2000\times2500 = \qty{5.00e6}{kg.m^{-2}.s^{-1}}$, folhelho
$2300\times2300 = \qty{5.29e6}{kg.m^{-2}.s^{-1}}$, calcário $2600\times4500 =
\qty{11.70e6}{kg.m^{-2}.s^{-1}}$. Giz--folhelho: $R = (5.29 - 5.00)/(5.29 + 5.00)
= 0.29/10.29 = 0.0282$; folhelho--calcário: $R = (11.70 - 5.29)/(11.70 + 5.29) =
6.41/16.99 = 0.377$. A energia refletida é $R^2 = 0.142$: o topo do
calcário devolve \qty{14.2}{\percent} da energia incidente, enquanto a
interface giz--folhelho devolve menos de \qty{0.1}{\percent}.
\end{solution}

\begin{solution}{exo:b3:exploration-seismology:2}
$t_0 = 2\times800/2000 = \qty{0.800}{s}$; $t = \sqrt{0.800^2 + (1000/2000)^2} =
\sqrt{0.89} = \qty{0.9434}{s}$; o sobretempo é \qty{0.1434}{s}.
A aproximação dá $1000^2/(2\times2000^2\times0.800) = \qty{0.1563}{s}$, alta
demais em \qty{9}{\percent}: para um afastamento maior que a profundidade, $x \ll vt_0$
deixa de valer.
\end{solution}

\begin{solution}{exo:b3:exploration-seismology:3}
Os pontos médios se deslocam pela metade do intervalo entre receptores: o intervalo entre pontos médios comuns é
$12.5/2 = \qty{6.25}{m}$. A multiplicidade é $240\times12.5/(2\times25) = 60$. Uma multiplicidade de
120 exige um disparo a cada $240\times12.5/(2\times120) = \qty{12.5}{m}$.
\end{solution}

\begin{solution}{exo:b3:exploration-seismology:4}
$\lambda = v/f$: \qty{100}{m} no arenito, $\lambda/4 = \qty{25}{m}$;
\qty{150}{m} no calcário, $\lambda/4 = \qty{37.5}{m}$. Com
$\lambda = 2800/30 = \qty{93.3}{m}$ e $z = \qty{1000}{m}$,
$r_F \approx \sqrt{\lambda z/2} = \sqrt{46\,667} = \qty{216}{m}$ (\qty{217}{m}
com o termo $\lambda^2/16$).
\end{solution}

\begin{solution}{exo:b3:exploration-seismology:5}
Pela \omterm{thm:b3:exploration-seismology:dix}{equação de Dix}: $v^2 = (2800^2\times0.700 - 2400^2\times0.500)/0.200 = (5\,488\,000 -
2\,880\,000)/0.200 = \qty{1.304e7}{m^2.s^{-2}}$, logo $v = \qty{3611}{m.s^{-1}}$.
A camada é atravessada em \qty{0.100}{s} no sentido de ida: $h = 3611\times0.100 =
\qty{361}{m}$.
\end{solution}

\begin{solution}{exo:b3:exploration-seismology:6}
(a) Seja $d$ a distância do ponto médio M ao refletor, de modo que
$t_0 = 2d/v$. O raio refletido da fonte S ao receptor G tem o mesmo
comprimento que a reta que liga a imagem S$'$ de S no refletor a G. Coloca-se M
na origem, com a linha ao longo do mergulho; S$'$ fica à distância $2d_S$ de S ao longo
da normal, com $d_S = d - (x/2)\sin\varphi$. O quadrado de S$'$G é
$x^2 + 4d_S^2 + 4d_Sx\sin\varphi = x^2 + 4d^2 - x^2\sin^2\varphi$, com todos os termos
em $dx$ se cancelando, logo $v^2t^2 = 4d^2 + x^2\cos^2\varphi$ e
$t^2 = t_0^2 + x^2\cos^2\varphi/v^2$. (b) O agrupamento segue uma hipérbole de velocidade
$V = v/\cos\varphi = 2000/\cos\ang{12} = \qty{2045}{m.s^{-1}}$: um mergulho faz com que a
\omterm{def:b3:exploration-seismology:velocities}{velocidade de empilhamento} fique alta demais, aqui em \qty{2.2}{\percent}.
\end{solution}

\begin{solution}{exo:b3:exploration-seismology:7}
Melhora-se por $\sqrt{60} = 7.75$. Dobrar a razão sinal-ruído exige quatro vezes
mais traços: uma multiplicidade de 240.
\end{solution}

\begin{solution}{exo:b3:exploration-seismology:8}
A profundidade plotada $vt/2$ varia em $\qty{1.5}{km.s^{-1}}\times\qty{0.40}{s}
= \qty{0.60}{km}$ por quilômetro: $\tan\theta_a = 0.60$, $\theta_a = \ang{31.0}$.
Pelo teorema do afastamento nulo $\sin\theta = 0.60$: $\theta = \ang{36.9}$.
\end{solution}

\begin{solution}{exo:b3:exploration-seismology:9}
Cada impulso coloca uma cópia de $w$ centrada nele: $0.3\,w$ nas amostras 0--2 dá
$(-0.15, 0.3, -0.15)$, e $-0.2\,w$ nas amostras 2--4 dá $(0.1, -0.2, 0.1)$.
A soma delas é $(0, -0.15, 0.3, -0.05, -0.2, 0.1, 0)$, escrita de uma amostra
antes da série até uma depois dela. As reflexões têm picos nas amostras 1 e 3, com
seus sinais verdadeiros; elas se sobrepõem na amostra 2, onde $-0.15 + 0.1 = -0.05$: uma
pequena interferência, mas dois picos distintos.
\end{solution}

\begin{solution}{exo:b3:exploration-seismology:10}
$\arcsin(1800/2900) = \ang{38.37}$ e $\arcsin(1800/4400) = \ang{24.15}$, logo
$i_c = (38.37 + 24.15)/2 = \ang{31.26}$ e $\theta = (38.37 - 24.15)/2 =
\ang{7.11}$; $v_2 = 1800/\sin\ang{31.26} = \qty{3469}{m.s^{-1}}$. A
média harmônica $2\times2900\times4400/(2900 + 4400) = \qty{3496}{m.s^{-1}}$ é
\qty{0.8}{\percent} maior; multiplicada por $\cos\theta = 0.9923$ dá
\qty{3469}{m.s^{-1}}, a aproximação usual para mergulhos pequenos
$v_2 \approx \cos\theta\,2v_dv_u/(v_d + v_u)$.
\end{solution}

\begin{solution}{exo:b3:exploration-seismology:11}
(a) O traço composto é $R[w(t) - w(t - \Delta t)] \approx R\,\Delta t\,w'(t)$ pela
expansão de Taylor de primeira ordem. (b) Sua amplitude é proporcional a
$\Delta t$, logo uma camada mais fina que a \omterm{def:b3:exploration-seismology:resolution}{espessura de sintonia} ainda devolve um sinal,
com a forma de $w'$, mas seu topo e sua base não podem ser identificados como dois eventos: ela
é detectada, não resolvida. Se $R$ e a ondaleta forem conhecidos, por exemplo a partir de um
poço, a amplitude dá $\Delta t$ e, portanto, a espessura. (c) A amplitude é
proporcional à espessura: a razão é $5/10 = 1/2$.
\end{solution}

\begin{solution}{exo:b3:exploration-seismology:12}
(a) $t(x) = (2/v)\sqrt{z^2 + x^2}$: $t(0) = 2\times600/3000 = \qty{0.400}{s}$ e
$t(300) = (2/3000)\sqrt{600^2 + 300^2} = \qty{0.4472}{s}$. (b) $t^2 = t_0^2 +
x^2/(v/2)^2$, a forma do \omterm{thm:b3:exploration-seismology:nmo}{sobretempo normal} com a velocidade $v/2$. Uma reflexão com o mesmo tempo de
ápice, em uma seção de afastamento nulo, não mostraria sobretempo algum; comparada com uma
hipérbole de \omterm{thm:b3:exploration-seismology:nmo}{sobretempo normal} com o mesmo $t_0$ em um agrupamento, o sobretempo da difração
$x^2/(2(v/2)^2t_0)$ é quatro vezes maior, porque o tempo depende da
distância ao espalhador duas vezes, na ida e na volta. (c) A migração soma a
seção ao longo desta hipérbole e põe a energia no seu ápice: a difração
colapsa para um ponto em $x = 0$, $z = \qty{600}{m}$.
\end{solution}

\begin{solution}{pb:b3:exploration-seismology:1}
\textbf{1.} De O, $400/(486.1 - 425.2) = 400/0.0609 = \qty{6568}{m.s^{-1}}$;
de L, $400/(486.1 - 372.4) = 400/0.1137 = \qty{3518}{m.s^{-1}}$. A onda de
cabeça vinda de O parece rápida: O é a extremidade no sentido do mergulho, e seus disparos percorrem o sentido contrário ao mergulho; o
refrator sobe de O para L.

\textbf{2.} $\arcsin(2400/3518) = \ang{43.02}$, $\arcsin(2400/6568) =
\ang{21.43}$: $i_c = \ang{32.22}$, $\theta = \ang{10.79}$.

\textbf{3.} $v_2 = 2400/\sin\ang{32.22} = \qty{4501}{m.s^{-1}}$: dentro da
faixa típica do calcário, o topo do calcário jurássico.

\textbf{4.} Os interceptos são $t - x/v_{\text{ap}}$ em \qty{1200}{m}:
$t_{i,O} = 0.4861 - 1200/6568 = \qty{0.3034}{s}$ e $t_{i,L} = 0.4861 -
1200/3518 = \qty{0.1450}{s}$. Com $\cos i_c = 0.8460$, $h = v_1t_i/(2\cos
i_c)$: $h_O = 2400\times0.3034/1.692 = \qty{430.4}{m}$ e $h_L =
\qty{205.7}{m}$.

\textbf{5.} Dividindo por $\cos\theta = 0.9823$: \qty{438.1}{m} sob O e
\qty{209.4}{m} sob L. Então $(438.1 - 209.4)/1200 = 0.1906 =
\tan\ang{10.79}$, o mergulho da questão 2.

\textbf{6.} O folhelho não produz onda de cabeça: ele é invisível, e o levantamento
o contabiliza junto com o giz. Aqui $v_1 = \qty{2.40}{km.s^{-1}}$ foi dada como a
média dos dois; se a velocidade da onda direta do giz, \qty{2.50}{km.s^{-1}},
tivesse sido usada para todo o pacote sobrejacente, todas as profundidades sairiam grandes demais.

\textbf{7.} $V = 600/\sqrt{0.5437^2 - 0.4878^2} = 600/0.2401 =
\qty{2499}{m.s^{-1}}$ e $V = 600/\sqrt{0.6240^2 - 0.5748^2} =
\qty{2471}{m.s^{-1}}$. O segundo valor combina o giz com o folhelho mais lento
logo abaixo dele: a média cai.

\textbf{8.} \qty{2811}{m.s^{-1}} e \qty{2829}{m.s^{-1}}.

\textbf{9.} Camada 1: \qty{2499}{m.s^{-1}}. Camada 2: $v^2 = (2471^2\times0.5748
- 2499^2\times0.4878)/0.0870$, $v = \qty{2308}{m.s^{-1}}$. Camada 3:
$v^2 = (2811^2\times0.6548 - 2471^2\times0.5748)/0.0800$,
$v = \qty{4561}{m.s^{-1}}$. Camada 4: $v^2 = (2829^2\times0.7348 -
2811^2\times0.6548)/0.0800$, $v = \qty{2972}{m.s^{-1}}$.

\textbf{10.} Giz (\qty{2499}{m.s^{-1}}), folhelho (\qty{2308}{m.s^{-1}}),
calcário (\qty{4561}{m.s^{-1}}), arenito (\qty{2972}{m.s^{-1}}), cada um dentro
da faixa publicada. As reflexões são a interface giz--folhelho, o topo
do calcário, a interface calcário--arenito e a discordância no
topo do embasamento de xisto.

\textbf{11.} Giz e folhelho têm quase a mesma impedância: $R = 0.028$
(exercício 1), contra 0.38 para o topo do calcário.

\textbf{12.} A hipérbole de velocidade $V_{\text{mq}}$ vale apenas até a segunda ordem
no afastamento; em afastamentos de \qty{2}{km}, o dobro da profundidade do
refletor mais profundo, os tempos se afastam dela, e uma hipérbole ajustada ao
espalhamento inteiro daria velocidades erradas, e os intervalos de Dix, ainda piores. Os
afastamentos longos também alongam as ondaletas ao serem corrigidos.

\textbf{13.} Tempos de ida e volta nas camadas: \qty{0.4878}{s}, \qty{0.0870}{s},
\qty{0.0800}{s}, \qty{0.0800}{s}. Espessuras $v\,\Delta t/2$: \qty{609.5}{m},
\qty{100.4}{m}, \qty{182.4}{m}, \qty{118.9}{m}.

\textbf{14.} Profundidades abaixo do solo: \qty{609.5}{m}, \qty{709.9}{m},
\qty{892.3}{m}, \qty{1011.2}{m}. O folhelho, o calcário e o arenito saem
dentro de \qty{0.4}{m}, \qty{2.4}{m} e \qty{1.1}{m} de suas espessuras, o
erro crescendo no calcário rápido e delgado, onde um pequeno erro de tempo custa
mais.

\textbf{15.} $\bar v = 2\times1011.2/0.7348 = \qty{2752}{m.s^{-1}}$. A \omterm{def:b3:exploration-seismology:velocities}{velocidade média quadrática}
pondera cada camada por $v^2$ e nunca é menor que a média ponderada pelo
tempo: $2829\times0.7348/2 = \qty{1039}{m}$, \qty{28}{m} fundo demais.

\textbf{16.} $\lambda = 2308/30 = \qty{76.9}{m}$, $\lambda/4 = \qty{19.2}{m}$;
A estimativa de Ricker é $2308/(5.2\times30) = \qty{14.8}{m}$. O folhelho tem
\qty{100}{m} de espessura, cinco vezes sua \omterm{def:b3:exploration-seismology:resolution}{espessura de sintonia}: seu topo e sua base são
resolvidos. Ele entraria em sintonia, pela regra do quarto de comprimento de onda, em $f = v/(4b) =
2308/400 = \qty{5.77}{Hz}$.

\textbf{17.} $\lambda = 2829/30 = \qty{94.3}{m}$, $z = \qty{1011}{m}$:
$r_F \approx \sqrt{94.3\times1011/2} = \qty{218}{m}$ (\qty{220}{m} com o
termo $\lambda^2/16$).

\textbf{18.} $t(x)^2 = t_0^2 + (2x/v)^2$, logo $v = 2x/\sqrt{t^2 - t_0^2} =
400/\sqrt{0.5732^2 - 0.550^2} = 400/0.1614 = \qty{2478}{m.s^{-1}}$.

\textbf{19.} $z = vt_0/2 = 2478\times0.550/2 = \qty{681}{m}$.

\textbf{20.} A profundidade plotada $vt/2$ varia $1.24\times0.167 = 0.207$ por
unidade de comprimento: $\theta_a = \arctan0.207 = \ang{11.70}$. O mergulho real tem
$\sin\theta = 0.207$: $\theta = \ang{11.95}$, o mergulho do flanco.

\textbf{21.} $d = 2480\times0.655/2 = \qty{812}{m}$; a migração desloca a
reflexão $d\sin\theta = 812\times0.207 = \qty{168}{m}$ no sentido contrário ao mergulho, em direção à
crista do anticlinal, e sobe de \qty{812}{m} para $d\cos\theta =
\qty{795}{m}$.

\textbf{22.} $2500\times0.160/2 = \qty{200}{m}$.

\textbf{23.} A partir das questões 13 e 14, a discordância do embasamento fica
\textbf{\qty[text-series-to-math=true]{1011}{m}} abaixo do solo no PMC S.
\end{solution}
